\section{多项式函数的均匀分布}

本节的目的在于证明以下的命题：
设$\{F_n\}_{n\in\mb{N}}$是斐波那契数列，$F_{n+2}=F_{n+1}+F_n,F_0=0,F_1=1$，则对于几乎所有的实数$\alpha$，$\{F_n\cdot\alpha\mod 1\}_{n\in\mb{N}}$在$[0,1)$中均匀分布，这等价于$\forall\, f\in C(S^1)(S^1\cong [0,1),S^1=\,\text{圆周}\,)$，$\frac{1}{N}\sum_{n=0}^{N-1}{f\left(F_n\cdot\alpha\right)}\to\int_{S^1}{f}\d x$.\par
其中要使用的最重要的是下列的定理：
\begin{theorem}{紧致度量空间上连续函数空间的可分性}
    设$X$是一个紧致度量空间，$C(X)=\{X\,\text{上的连续函数}\}$，定义$C(X)$上的范数为$\|f\|=\sup_{x\in X}|f(x)|$，定义$d(f,g)=\|f-g\|$，则$(C(X),\|\cdot\|)$是一个可分的度量空间，即存在$\{f_k\}_{k\in\mb{N}}\subseteq C(X)$，使得$\{f_k\}_{k\in\mb{N}}$在$C(X)$中稠密.
\end{theorem}
考虑矩阵$A=\begin{pmatrix}
    0 & 1 \\
    1 & 1
\end{pmatrix}$和$T_A:X_2\to X_2$，则$|\lambda I-A|=\begin{vmatrix}
    \lambda & -1 \\
    -1 & \lambda-1
\end{vmatrix}=\lambda(\lambda-1)=\lambda^2-\lambda-1=0$，$A$的特征值均不是单位根，因此$T_A$是遍历映射.\par
由Birkhoff遍历定理，$\forall\, f\in L^1(X_2,m_\mathrm{Leb})$，$\frac{1}{N}\sum_{n=0}^{N-1}{f \left(T^n_A(x)\right)}\to \int_{X_2}{f}\d m_\mathrm{Leb}$，$\aew x+\mb{Z}^2\in X_2\xlongrightarrow{\text{提升到}}\aew x\in\mb{R}^2$.\par

根据上面的定理，$C(X_2)\subseteq L^1(X_2,m_\mathrm{Leb})$是可分空间，设$\{f_k\}_{k\in\mb{N}}\subseteq C(X_2)$是其中的一个可数稠密子集.\par
$\forall\, k\in\mb{N}$，存在$\mb{R}^2$中的全测集$Q_k$，使得$\forall\, x\in Q_k,\frac{1}{N}\sum_{n=0}^{N-1}{f_k \left(T^n_A(x)\right)}\to \int_X{f_k}\d m_\mathrm{Leb}$，取$Q=\bigcap_{k=1}^\infty{Q_k}$，则$Q$也是$\mb{R}^2$中的全测集.\par
$\forall\, f\in C(X_2),\forall\, \varepsilon>0$，由$\{f_k\}_{k\in\mb{N}}$稠密知$\exists\, k\in\mb{N},\st d(f_k,f)<\varepsilon \Leftrightarrow \sup_{x\in X_2}{|f(x)-f_k(x)|}<\varepsilon$.\par
$\forall\, x\in Q$，$\frac{1}{N}\sum_{n=0}^{N-1}{f_k \left(T^n_A(x)\right)}\to \int_{X_2}{f_k}\d m_\mathrm{Leb}$$\Rightarrow \varlimsup_{n\to\infty}{\left|\frac{1}{N}\sum_{n=0}^{N-1}{f \left(T^n_A(x)\right)}-\int_{X_2}{f}\d m_\mathrm{Leb}\right|}\leq 2\varepsilon $.\par
令$\varepsilon\to 0$，$\frac{1}{N}\sum_{n=0}^{N-1}{f \left(T^n_A(x)\right)}\to \int_{X_2}{f}\d m_\mathrm{Leb}$.\par
接下来研究$T^n_A(x)$.按照定义$T^{n+1}_A(x)=A^{n+1}x+\mb{Z}^2$.\par
其中$A^{n+1}=\underbrace{\begin{pmatrix}
    0 & 1\\
    1 & 1
\end{pmatrix}\begin{pmatrix}
    0 & 1\\
    1 & 1
\end{pmatrix}\cdots\begin{pmatrix}
    0 & 1\\
    1 & 1
\end{pmatrix}}_{n+1\,\text{个}}=\begin{pmatrix}
    q_{n-1} & q_{n} \\
    p_{n-1} & p_{n}
\end{pmatrix}$，其中$q_{n+1}=a_{n+1}q_n+q_{n-1}=q_n+q_{n-1},p_{n+1}=a_{n+1}p_n+p_{n-1}=p_n+p_{n-1}$满足斐波那契数列的递推公式.\par
当$n=0$时，$\begin{pmatrix}
    0 & 1 \\
    1 & 1 \\
\end{pmatrix}=\begin{pmatrix}
    q_{-1} & q_{0} \\
    p_{-1} & p_{0}
\end{pmatrix}\Rightarrow q_0=1,q_{-1}=0,p_0=1,p_{-1}=1$.\par
因此对于$\{p_n\}_{n\in\mb{N}}$，有$p_n=F_{n+2}$.\par
$\forall\, x\in Q,x=\begin{pmatrix}
    \alpha \\
    \beta
\end{pmatrix}\in\mb{R}^2$，$A^{n+1}x=\begin{pmatrix}
    q_{n-1} & q_{n} \\
    p_{n-1} & p_{n}
\end{pmatrix}\begin{pmatrix}
    \alpha \\
    \beta
\end{pmatrix}=\begin{pmatrix}
    \alpha q_{n-1}+\beta q_{n} \\
    \alpha p_{n-1}+\beta p_{n}
\end{pmatrix}$.\par
取$g,h\in C(S^1)$，则$f(u,v)=g(u)\cdot h(v)\in C(X_2)$.\par
则$\forall\, x\in\begin{pmatrix}
    \alpha \\
    \beta
\end{pmatrix}\in Q,\frac{1}{N}\sum_{n=0}^{N-1}{f g(\alpha q_{n-1}+\beta q_{n})h(\alpha p_{n-1}+\beta p_{n})}\to (\int_{S_1}{g}\d x)(\int_{S^1}{h}\d x)$.\par
令$g=\bf{1}$，则$\forall\, h\in C(S^1),\forall\, x=\begin{pmatrix}
    \alpha \\
    \beta
\end{pmatrix}\in Q,\frac{1}{N}\sum_{n=0}^{N-1}{h(\alpha p_{n-1}+\beta p_{n})}\to \int_{S^1}{h}\d x$.\par
$\Rightarrow \forall\,\begin{pmatrix}
    \alpha \\
    \beta
\end{pmatrix}\in Q,\{\alpha p_{n-1}+\beta p_{n}\}_{n\in\mb{N}}$在$[0,1)$中均匀分布.\par
由于$p_n=F_{n+2}$，所以$\{\alpha F_{n}+\beta F_{n+1}\}_{n\in\mb{N}}$也在$[0,1)$中均匀分布.\par
我们知道$F_n$的通项公式为$F_n=c(\lambda_1^n+\lambda_2^n)$，其中$|\lambda_1|>1,|\lambda_2|<1$.\par
那么$\setjot[-4]\begin{aligned}[t]
    \alpha F_{n}+\beta F_{n+1} &=\alpha c(\lambda_1^{n}+\lambda_2^{n})+\beta c(\lambda_1^{n+1}+\lambda_2^{n+1})=\lambda_1^n(c\beta\lambda_1+c\alpha)+(\lambda_2^n\cdot c\alpha+c\beta\cdot \lambda_2^{n+1}) \\
    &= F_n(\beta\lambda_1+\alpha)+\tilde{c}\lambda_2^n(\tilde{c}\,\text{为某个常数}) \\
    &=\tilde{\alpha}F_n+\tilde{c}\lambda_2^n
\end{aligned}$\par
因此$\frac{1}{N}\sum_{n=0}^{N-1}{h(\tilde{\alpha}F_n)}=\frac{1}{N}\sum_{n=0}^{N-1}{h(\alpha F_{n}+\beta F_{n+1}-\tilde{c}\lambda_2^n)}$\par
由于$|\lambda_2|<1$，故$|\lambda_2|^n\to 0$.\par
又$h$是连续函数，因此$\forall\,\varepsilon>0$，$\exists\, M\in\mb{N}$，使得$\forall\, n> M$，$\left|h(\tilde{\alpha}F_n)-h\left(\alpha F_n+\beta F_{n+1}\right)\right|<\varepsilon$.\par
所以$\setjot[-2]\begin{aligned}[t]
    &\quad\left|\frac{1}{N}\sum_{n=0}^{N-1}h\left(\tilde{\alpha}F_n\right)-\frac{1}{N}\sum_{n=0}^{N-1}h\left(\alpha F_n+\beta F_{n+1}\right)\right| \\
    &\leq\frac{1}{N}\sum_{n=0}^M{\left|h\left(\tilde{\alpha}F_n\right)-h\left(\alpha F_n+\beta F_{n+1}\right)\right|}+\frac{1}{N}\sum_{n=M+1}^N{\left|h\left(\tilde{\alpha}F_n\right)-h\left(\alpha F_n+\beta F_{n+1}\right)\right|} \\
    &\leq\frac{2M\|h\|}{N}+\varepsilon
\end{aligned} $\par
令$N\to\infty,\varepsilon\to 0$，则$\left|\frac{1}{N}\sum_{n=0}^{N-1}h\left(\tilde{\alpha}F_n\right)-\frac{1}{N}\sum_{n=0}^{N-1}h\left(\alpha F_n+\beta F_{n+1}\right)\right|\to 0$.\par
因此$\frac{1}{N}\sum_{n=0}^{N-1}{h(\tilde{\alpha}F_n)}\to\int_{S^1}h\d x$，故$\{\tilde{\alpha}F_n\}_{n\in\mb{N}}$在$[0,1)$中均匀分布.\par
最后只需证明上述$\tilde{\alpha}$的全体构成一个全测集，即$m_\mathrm{Leb}\left(\left\{\tilde{\alpha}\left|\tilde{\alpha}=\alpha\lambda_1+\beta,\begin{pmatrix}
    \alpha \\
    \beta
\end{pmatrix}\in Q\right.\right\}\right)=1$.

\begin{proof}
$Q\subseteq \mb{R}^2$为全测集，由Fubini定理，存在$A\subseteq\mb{R}$的全测集，使得$\forall\,\alpha\in A$，其“垂直切片”$Q_\alpha:=\{\beta\in\mb{R}:(\alpha,\beta)\in Q\}$
在$\mb{R}$中也是全测集.\par
固定任意$\alpha_0\in A$，令$E:=\{\tilde{\alpha}\in\mb{R}\mid \tilde{\alpha}=\alpha_0\lambda_1+\beta,\ \beta\in Q_{\alpha_0}\}=\alpha_0\lambda_1+Q_{\alpha_0}$. \par
由Lebesgue测度的平移不变性，$E$在$\mb{R}$中是全测集.\par
又$E\subseteq\{\tilde{\alpha}\in\mb{R}\mid \tilde{\alpha}=\alpha\lambda_1+\beta,(\alpha,\beta)\in Q\}$，
故$\{\tilde{\alpha}\mid \tilde{\alpha}=\alpha\lambda_1+\beta,(\alpha,\beta)\in Q\}$也是$\mb{R}$中全测集.\par
\end{proof}
至此，我们就证明了本节一开始提出的命题：
\begin{proposition}
    设$\{F_n\}_{n\in\mb{N}}$为斐波那契数列，则对于几乎所有的实数$\alpha$，$\{F_n\cdot\alpha\mod 1\}_{n\in\mb{N}}$在$[0,1)$中均匀分布；等价地，$\forall\, f\in C(S^1)$，$\frac{1}{N}\sum_{n=0}^{N-1}{f\left(F_n\cdot\alpha\right)}\to\int_{S^1}{f}\d x$.
\end{proposition}

本节的第二个例子是下面这个定理：
\begin{theorem}
    设$\a\in\mb{R}\setminus\mb{Q}$，$p(x)=\alpha x^n+a_{n-1}x^{n-1}+\cdots+a_1x+a_0(a_i\in\mb{R})$，则$\{p(n)\mod 1\}_{n\in\mb{N}}$在$[0,1)$中均匀分布.
\end{theorem}
根据均匀分布的等价条件，我们实际上只需要证明$\forall\, k\in\mb{Z}\setminus\{0\},\frac{1}{N}\sum_{n=0}^{N-1}{\cos{2k\pi p(n)}}\to 0$且$\frac{1}{N}\sum_{n=0}^{N-1}{\sin{2k\pi p(n)}}\to 0$.\par
这等价于证明$\frac{1}{N}\sum_{n=0}^{N-1}{\eu^{2k\pi \iu p(n)}}\to 0(\forall\, k\in\mb{Z}\setminus\{0\})$.\par
对于上面的定理，如果$p(n)=\alpha n+a_0$，它在$[0,1)$上均匀分布，因而$\frac{1}{N}\sum_{n=0}^{N-1}{\eu^{2k\pi \iu(\alpha n+a_0)}}\to 0$.为了将其余次数的多项式化归到这一情形，我们需要下述引理.\par
\begin{lemma}
    设$\{a_n\}_{n\in\mb{N}}$是$\mb{C}$中的有界序列，且满足$\forall\, h\in\mb{N}_{>0}$，$\frac{1}{N}\sum_{n=0}^{N-1}{\overline{a_{n+h}}a_n}\to 0$，那么$\frac{1}{N}\sum_{n=0}^{N-1}{a_n}\to 0$.
\end{lemma}
\begin{proof}
    记$b_n=\frac{1}{H}\sum_{h=0}^{H-1}{a_{n+h}}(\forall\, H\in\mb{N})$，则$b_n$实际上就是从$a_n$开始的长度为$H$的$a_i$序列的平均值.\par
    $\setjot[-2]\begin{aligned}[t]
        \left|\frac{1}{N}\sum_{n=0}^{N-1}a_n-\frac{1}{N}\sum_{n=0}^{H-1}b_n\right|&\leq \frac{1}{N}\left(\sum_{n=0}^{N-1}|a_1|+\sum_{n=N-H}^{N-1}|a_n|\right)+\frac{1}{N}\left(\sum_{n=0}^{H-1}|a_1|+\sum_{n=N}^{N+H-1}|a_n|\right) \\
        &\leq \frac{4H}{N}M\;(|a_i|\leq M)
    \end{aligned}$\par
    对于固定的$H\in\mb{N}$，$\lim_{n\to\infty}\left|\frac{1}{N}\sum_{n=0}^{H-1}a_n-\frac{1}{N}\sum_{n=0}^{N-1}{b_n}\right|=0$.\par
    由均值不等式可知$\left|\frac{1}{N}\sum_{n=0}^{N-1}b_n\right|\leq \sqrt{\frac{1}{N}\sum_{n=0}^{N-1}|b_n|^2}$.\par
    而$\setjot[-2]\begin{aligned}[t]
        \frac{1}{N}\sum_{n=0}^{N-1}|b_n|^2 &=\frac{1}{N}\sum_{n=0}^{N-1}\left|\frac{1}{N}\sum_{h=0}^{H-1}a_{n+h}\right|^2=\frac{1}{N}\frac{1}{H^2}\sum_{n=0}^{N-1}\left(\sum_{i=0}^{H-1}\overline{a_{n+i}}\right)\left(\sum_{j=0}^{H-1}{a_{n+j}}\right) \\
        &=\frac{1}{N}\frac{1}{H^2}\sum_{n=0}^{N-1}\sum_{i=0}^{H-1}\sum_{j=0}^{H-1}\overline{a_{n+i}}a_{n+j}=\frac{1}{H^2}\sum_{0\leq i,j\leq H-1}\left(\frac{1}{N}\sum_{n=0}^{N-1}{\overline{a_{n+1}}a_{n+j}}\right) \\
        &=\frac{1}{H^2}\sum_{0\leq i=j\leq H-1}\left(\frac{1}{N}\sum_{n=0}^{N-1}{\overline{a_{n+1}}a_{n+j}}\right)+\frac{1}{H^2}\sum_{\substack{0\leq i,j\leq H-1\\ i\neq j}}\left(\frac{1}{N}\sum_{n=0}^{N-1}{\overline{a_{n+1}}a_{n+j}}\right)
    \end{aligned}$\par
    固定$H\in\mb{N}$，则$\varlimsup_{N\to\infty}{\frac{1}{N}\sum_{n=0}^{N-1}{|b_n|^2}}\leq\frac{M}{H}\quad \left(|a_i|\leq M,\forall\, i\in\mb{N}\right)$.\par
    $\Rightarrow \varlimsup_{N\to\infty}\left|\frac{1}{N}\sum_{n=0}^{N-1}{b_n}\right|\geq\frac{M}{\sqrt{H}}$.\par
    令$H\to\infty$，得$\setjot[-2]\begin{aligned}[t]
        \varlimsup_{N\to\infty}\left|\frac{1}{N}\sum_{n=0}^{N-1}a_n\right|&\leq\varlimsup_{N\to\infty}\left|\frac{1}{N}\sum_{n=0}^{N-1}{a_n}-\frac{1}{N}\sum_{n=0}^{N-1}b_n\right|+\varlimsup_{N\to\infty}\left|\frac{1}{N}\sum_{n=0}^{N-1}b_n\right| \\
        &\leq 0+\frac{M}{\sqrt{H}} \to 0\;(H\to\infty)
    \end{aligned}$\par
\end{proof}

借助这个引理，我们就可以证明上述的定理，先以三次多项式为例.
\begin{instance}
    若$p(n)=\alpha n^3+n$，需要证明$\frac{1}{N}\sum_{n=0}^{N-1}{\eu^{2k\pi \iu p(n)}}\to 0\;(\forall\, k\in\mb{Z}\setminus\{0\})$.\par
    设$a_n=\eu^{2k\pi \iu p(n)}$，即证$\frac{1}{N}\sum_{n=0}^{N-1}{\overline{a_{n+h}}a_n}\to 0\;(\forall\, h\in\mb{N}_{>0})$. \par
    其中$\overline{a_{n+h}}a_n=\eu^{-2k\pi \iu p(n+h)}\eu^{2k\pi \iu p(n)}=\eu^{-2k\pi \iu \left(\alpha(3n^2h+3nh^2+h^3)+h\right)}$.\par
    记$b_n=\eu^{-2k\pi \iu \left(\alpha(3n^2h+3nh^2+h^3)+h\right)}$，需要证明$\frac{1}{N}\sum_{n=0}^{N-1}{b_n}\to 0\;(\forall\, h\in\mb{N}_{>0})$.\par
    $\overline{b_{n+s}}b_n=\eu^{-2k\pi\iu \left(\alpha(3h(2ns+s^2)+3sh^2)\right)}$，这是关于$n$的一次多项式，因此$\frac{1}{N}\sum_{n=0}^{N-1}{\overline{b_{n+s}}b_n}\to 0\;(\forall\, s\in\mb{N}_{>0})$.\par
    因此$\frac{1}{N}\sum_{n=0}^{N-1}b_n=\frac{1}{N}\sum_{n=0}^{N-1}{\overline{a_{n+h}a_n}}\to 0\;(\forall\,h\in\mb{N}_{>0})$. \par
    $\Rightarrow \frac{1}{N}\sum_{n=0}^{N-1}{a_n}\to 0\;(\forall\, k\in\mb{Z}\setminus\{0\})$.
\end{instance}
类似上面的例子，对于一般的多项式$p(n)$，只需要对$\mathrm{deg}\, p$用数学归纳法即可.另外，上述定理实际上可以加强为：
\begin{theorem}{}
    $p(x)=a_nx^n+a_{n-1}x^{n-1}+\cdots+a_1x+a_0$，若$a_i(1\leq i\leq n)$不全为有理数，则$\{p(n)\mod 1\}_{n\in\mb{N}}$在$[0,1)$中均匀分布.
\end{theorem}
\begin{proof}
    我们对多项式的次数$n$使用第二数学归纳法.\par
    当$n=1$时，$p(x)=a_1x+a_0$，其中$a_1$是无理数，成立.\par
    假设定理对所有次数小于$n$的多项式成立，考虑$n$次多项式$p(x)=a_nx^n+\cdots+a_0$.\par
    设$m$是使得$a_m$为无理数的最大下标($1\leq m\leq n$).\par
    要证$\frac{1}{J}\sum_{j=0}^{J-1}{\eu^{2k\pi \iu p(j)}}\to 0\;(\forall\, k\in\mb{Z}\setminus\{0\})$，设$b_j=\eu^{2k\pi \iu p(j)}$，即证$\frac{1}{J}\sum_{j=0}^{J-1}{\overline{b_{j+h}}b_j}\to 0\;(\forall\, h\in\mb{N}_{>0})$，其中$\overline{b_{j+h}}b_j=\eu^{-2k\pi\iu(p(j+h)-p(j))}$. \par
    记$q_h(j)=p(j+h)-p(j)$，$q_h(j)=\sum_{k=1}^n a_k ((j+h)^k - j^k)$.\par
    考虑$q_h(j)$中$j^{m-1}$的系数，它来自于$q_h(x)$中$k\geq m$的项展开式中.\par
    由二项式定理$a_k((j+h)^k - j^k) = a_k\left(kj^{k-1}h + \binom{k}{2}j^{k-2}h^2 + \cdots\right)$.\par
    因此，$q_h(j)$中$j^{m-1}$的系数为$c_{m-1} = mha_m + \sum_{k=m+1}^n a_k \cdot \binom{k}{k-m+1}h^{k-m+1}$.\par
    由$\forall\,k>m$，$a_k\in\mb{Q}$且$h\in\mb{N}_{>0}$，可知$\sum_{k=m+1}^n a_k \cdot \binom{k}{k-m+1}h^{k-m+1}$ 是有理数.\par
    而$a_m$是无理数，所以$mha_m$是无理数.故$q_h(j)$的$j^{m-1}$次项系数$c_{m-1}$是无理数.\par
    $q_h(j)$是次数为$n-1$的多项式且$j^{m-1}$的系数为无理数.根据假设，对于$\forall\, h>0$，$q_h(j) \mod 1$在$[0,1)$中均匀分布，即$\frac{1}{J}\sum_{j=0}^{J-1}{\overline{b_{j+h}}b_j}\to 0\;(\forall\, h\in\mb{N}_{>0})$. \par
    因此，$\frac{1}{J}\sum_{j=0}^{J-1}{\eu^{2k\pi \iu p(j)}}\to 0\;(\forall\, k\in\mb{Z}\setminus\{0\})$，$p(x)\mod 1$在$[0,1)$中均匀分布.\par 
    由归纳法可知上述命题对$n\in\mb{N}_{n>0}$均成立.
\end{proof}


\section{回复定理}

下面学习一些回复定理并使用回复定理回答一些问题.\par
\begin{theorem}{Poincare回复定理}
    设$(X,\mcal{B},m)$为概率测度空间，$T:X\to X$是保测映射，$A\in\mcal{B},m(A)>0$，那么存在$F\subseteq A,m(A\setminus F)=0$，使得$\forall\, x\in F$，从$x$出发的轨迹$\{x,Tx,T^2x,\cdots,T^nx,\cdots\}$无限多次回到$A$中(更进一步，回到$F$中）.
\end{theorem}
\begin{proof}
    记$B_1=T^{-1}A\cup T^{-2}A\cup\cdots\cup T^{-n}A\cup\cdots$，即$\forall\, x\in B,\exists\, n\in\mb{N},T^nx\in A$.\par
    $B_2=T^{-2}A\cup T^{-3}A\cup\cdots\cup T^{-(n+1)}A\cup\cdots$，即$\forall\, x\in B_2,\exists\, n\in\mb{N},n\geq 2,T^nx\in A$.\par
    $\quad\vdots$\par
    $B_k=T^{-k}A\cup T^{-(k+1)}A\cup\cdots$，以此类推.\par
    令$B_0=A\cup T^{-1}A\cup\cdots\cup T^{-n}A\cup\cdots$.\par
    则$B_0\supseteq B_1\supseteq B_2\supseteq\cdots\supseteq B_k\supseteq\cdots$，且$T^{-1}B_k=B_{k+1},m(B_k)=m(B_{k+1})$.\par
    令$B_\infty=\bigcap_{i=0}^\infty{B_i}=\bigcap_{i=1}^\infty{B_i}$，那么$m(B_\infty)=m(B_k)$.\par
    则$\forall\, x\in B_\infty$，$x$出发的轨道无限次进入集合$A$.\par
    令$F=B_\infty\cap A$，由$m(B_\infty)=m(B_0),B_0=A\cup T^{-1}A\cup T^{-2}A\cup\cdots\supseteq A$，得$m(F)=m(A)$.\par
    下面证明$\forall\, x\in F$，$x$出发的轨道无限次进入集合$F$.\par
    $\forall\, x\in F$，存在子列$n_1<n_2<\cdots<n_k<\cdots$，使得$T^{n_1}x\in A,T^{n_2}x\in A,\cdots,T^{n_k}x\in A,\cdots$.\par
    对于$T^{n_1}x$，$T^{n_i}x=T^{n_i-n_1}\left(T^{n_1}x\right)\in A\;(\forall\, i\geq 2)$，因此$T^{n_1}x\in F$.\par
    类似地，$T^{n_2}x,T^{n_3}x,\cdots\in F$.
\end{proof}

下面的回复定理跳出了概率测度空间的范畴，它定义在度量空间上.
\begin{theorem}{Birkhoff回复定理}
    设$(X,d)$为紧致度量空间，$T:X\to X$是连续映射，存在$x\in X$，以及自然数子列$\{n_i\}_{i\in\mb{N}}$，使得$T^{n_i}x\to x(i\to\infty)$
\end{theorem}

\begin{definition}{$\omega-$极限集和$\alpha-$极限集}
    $\forall\, x\in X$，记$\omega(x)$为序列$Tx,T^2x,\cdots,T^nx,\cdots$所有极限点构成的集合，称为$x$的\textbf{$\bm{\omega}$--极限集}.\par
    如果$T$为同胚，则记$\alpha(x)$为序列$T^{-1}x,T^{-2}x,\cdots,T^{-n}x,\cdots$所有极限点构成的集合，称为$x$的\textbf{$\bm{\alpha}$--极限集}.
\end{definition}

\begin{proposition}{$\omega$--极限集的性质}
    \begin{enumerate}
        \item $\forall\, x\in X,\omega(x)\neq\varnothing$.
        \item $\omega(x)$是闭集(紧集).
        \item $\omega(x)$是$T-$不变的，即$T\omega(x)=\omega(x)$.
    \end{enumerate}
\end{proposition}
\begin{proof}
    $\forall\, y\in T\omega(x)$，$\exists\, z\in\omega(x)$，使得$y=Tz$.\par
    设$z=\lim_{i\to\infty}T^{n_i}x\;(n_1<n_2<\cdots<n_k<\cdots)\Rightarrow y=\lim_{i\to\infty}T^{n_i+1}x\Rightarrow y\in\omega(x)$.\par
    因此$T\omega(x)\subseteq \omega(x)$.\par
    $\forall\, y\in\omega(x)$，$\exists\, n_1<n_2<\cdots<n_k<\cdots$，使得$\lim_{i\to\infty}{T^{n_i}x}=y$，则$\lim_{i\to\infty}T \left(T^{n_i-1}(x)\right)=y$.\par
    由$X$的紧致性，取$\left\{T^{n_i-1}(x)\right\}$的收敛子列$\left\{T^{m_i}(x)\right\}$，并记$T^{m_i}(x)\to z$.\par
    那么由连续性$y=Tz\Rightarrow y\in T\omega(x)$.\par
    因此$\omega(x)\subseteq T\omega(x)$.
\end{proof}

\begin{definition}{极小系统}
    设$(X,d)$是紧致度量空间，$T:X\to X$是连续映射，如果$\forall\, x\in X,\left\{x,Tx,\cdots,T^nx,\cdots\right\}$在$X$中稠密，那么称$T:X\to X$为\textbf{极小系统}.\par
    特别地，若$T:X\to X$是同胚，如果$\forall\, x\in X, \left\{\cdots,T^{-2}x,T^{-1}x,x,Tx,T^2x,\cdots\right\}$在$X$中稠密，那么称$T:X\to X$为\textbf{极小系统}.
\end{definition}

\begin{lemma}{Zorn引理}
    设$\mcal{F}$是一个集合，$\prec $是$\mcal{F}$上的一个偏序关系.如果对于任意一个全序子集$A\subseteq\mcal{F}$，可以找到元素$x$使得$\forall\, y\in A,x<y$，那么$\mcal{F}$中存在极小元.
\end{lemma}

\begin{theorem}
    设$(X,d)$为紧致度量空间，$T:X\to X$是连续映射，那么存在$X$中的非空闭集$A$，使得$TA\subseteq A$，且$T:A\to A$是极小系统.
\end{theorem}
\begin{proof}
    设$\mcal{F}=\left\{A\subseteq X|A\,\text{非空闭集}\, ,TA\subseteq A\right\}$.\par
    定义偏序关系：$A<B\iff A\subseteq B(\forall\, A,B\in\mcal{F})$.\par
    设$\mcal{A}$是$\mcal{F}$中的全序子集，$\mcal{A}=\{A_\alpha\}_{\alpha\in I}$.\par
    定义$E=\bigcap_{\alpha\in I}{A_\alpha}$，下证$E\subseteq \mcal{F}$.\par
    \begin{enumerate}[leftmargin=4em,label=(\arabic*)]
        \item $E$是非空闭集.\par
        由$A_\alpha$闭，$E=\bigcap_{\alpha\in I}{A_\alpha}$为闭集.\par
        若$E=\varnothing$，取补得$X=\bigcup_{\alpha\in I}{A_\alpha}$为闭集.\par
        由于$X$ 紧致，因此存在$\{A_\alpha^c\}_{\alpha\in I}$ 的有限子覆盖$\left\{A_1^c,A_2^c,\cdots,A_k^c\right\}$ ，不妨$A_1<A_2<\cdots A_k$.\par
        则$X=A_1^c\cup A_2^c\cup \cdots\cup A_k^c$，取补得$\varnothing=A_1\cap A_2\cap \cdots A_k=A_1$，矛盾! \par
        因此$E$ 非空.
        \item $TE\subseteq E$.\par
        $TA_\alpha\subseteq A_\alpha(\forall\,\alpha\in I)\Rightarrow T \left(\bigcap_{\alpha\in I}{A_\alpha}\right)\subseteq \bigcap_{\alpha\in I}{A_\alpha}\Rightarrow TE\subseteq E$.
    \end{enumerate}\par
    因此，$E\subseteq \mcal{F}$.\par
    由Zorn引理，$\mcal{F}$中存在极小元，设为$A$.\par
    定义$T|_A:A\to A$，下证其为极小系统，即$\forall\, x\in A$，$\left\{x,Tx,\cdots,T^nx,\cdots\right\}$在$A$中稠密.\par
    设$S=\overline{\{x,Tx,T^2x,\cdots\}}$.\par
    由于$\forall\, i\in\mb{N}_{>0},T(T^i(x))=T^{i+1}x\in S\Rightarrow T \left(\overline{\{x,Tx,\cdots\}}\right)\subseteq S$，因此$TS\subseteq S$.\par
    所以$S\in\mcal{F}$，如果$S\neq A$，那么$S<A$，与$A$的极小性矛盾.\par
    故$S=\overline{\{x,Tx,T^2x,\cdots\}}=A\Rightarrow \{x,Tx,T^2x,\cdots\}$ 在$A$中稠密.
\end{proof}

基于该定理，我们就可以证明Birkhoff回复定理.
\begin{theorem*}{Birkhoff回复定理}
    设$(X,d)$为紧致度量空间，$T:X\to X$是连续映射，存在$x\in X$，以及自然数子列$\{n_i\}_{i\in\mb{N}}$，使得$T^{n_i}x\to x(i\to\infty)\Leftrightarrow x\in \omega(x)$.
\end{theorem*}
\begin{proof}
    不妨假设$T:X\to X$是极小系统(否则考虑$T|_A:A\to A$即可).\par
    $\forall\, x\in X,x$的$\omega$--极限集$\omega(x)$是闭集，且$T \left(\omega(x)\right)=\omega(x)$.\par
    由于$T$是极小系统，$\forall\, y\in X$，$\overline{\{y,Ty,T^2y,\cdots\}}=X$.\par
    取$y\in\omega(x)$，由$T$--不变性$\{y,Ty,T^2y,\cdots\}\subseteq \omega(x)$.\par
    又因$\omega(x)$是闭集，故$\overline{\{y,Ty,T^2y,\cdots\}}\subseteq\omega(x)$，即$X\subseteq\omega(x)$.\par
    因此$\omega(x)=X$，这意味着$x\in\omega(x)$.
\end{proof}

下面介绍上述回复定理的推广，上面的定理对于多个映射的情形类似可成立.\par
\begin{theorem}{多重Birkhoff回复定理}
    设$(X,d)$ 是紧致度量空间，$T:X\to X$ 同胚，$l\in\mb{N}$，那么$\exists\, x\in X$，以及自然数子列$n_1<n_2<\cdots<n_k<\cdots$，使得$T^{n_i}x\to x,\left(T^2\right)^{n_i}x\to x,\cdots,\left(T^l\right)^{n_i}x\to x$.
\end{theorem}

为了证明这个定理，需要一系列概念和引理.\par
\begin{definition}{回复集}
    设$(X,d)$ 是紧致度量空间，$T:X\to X$ 连续映射，$A\subseteq X$ 非空闭集，如果$\forall\,\varepsilon>0,\forall\, x\in A$ ，存在$y\in A$ 和$n\in\mb{N}_{>0}$ ，使得$d \left(T^ny,x\right)<\varepsilon$，则称$A$为\textbf{回复集}.
\end{definition}
可以注意到如果$A$ 只包含一个元素，那么这个元素就满足Birkhoff回复定理的结论，是一个回复点，因此回复集是回复点的推广.\par

下面一个引理指出回复集中存在点可以回到任意的“$\varepsilon$ 精度”，但这并不等同于回复集中存在回复点，因为其中的点和$\varepsilon$有关.\par 
\begin{lemma}{}{2.8}
    设$(X,d)$为紧致度量空间，$T:X\to X$是连续映射，如果$A\subseteq X$是回复集，那么$\forall\, \varepsilon>0,\exists\, z\in A$和$n\in\mb{N}_{>0}$ ，使得$d \left(T^{n}z,z\right)<\varepsilon$.
\end{lemma}
\begin{proof}
    设$z_0\in A$，$\forall\, \varepsilon>0$，由于$A$ 是回复集，存在$z_1\in A,n_1\in\mb{N}_{>0}$ ，使得$d \left(T^{n_1}z_1,z_0\right)<\varepsilon$.\par
    由于$T$ 是连续映射，$\exists\, \varepsilon_1>0$，使得$T^{n_1}\left(B(z_1,\varepsilon_1)\right)\subseteq B(z_0,\varepsilon)$.\par
    由于$A$ 是回复集，$\exists\, z_2\in A,n_2\in\mb{N}_{>0}$，使得$d \left(T^{n_2}z_2,z_1\right)<\varepsilon_1$.\par
    $\Rightarrow \exists\, \varepsilon_2>0$ ，使得$T^{n_2}\left(B(z_2,\varepsilon_2)\right)\subseteq B(z_1,\varepsilon_1)$.\par
    重复以上过程，可以构造点列$\{z_i\}_{i\in\mb{N}}$ 和序列$\{\varepsilon_i\}_{i\in\mb{N}}$ ，使得$T^{n_i}\left(B(z_i,\varepsilon_i)\right)\subseteq B(z_{i-1},\varepsilon_{i-1})(\forall\, i\in\mb{N}_{>0})$.\par
    由于$\{z_i\}_{i\in\mb{N}}\subseteq X$紧空间，因此存在收敛子列$z_{k_i}\to\omega\in X(i\to\infty)$.\par
    $\forall\,\varepsilon>0,\exists\, N>0,\st\forall\, i>N,d \left(z_{k_i},\omega\right)<\varepsilon$.\par
    则$\exists\, M_i\in\mb{N}$ ，使得$T^{M_i}{z_{k_i}}\in B \left(z_{k_{N+1}},\varepsilon_{k_{N+1}}\right)\subseteq B \left(z_{k_i},2\varepsilon+\varepsilon_{k_{N+1}}\right)$.\par
    $\Rightarrow d \left(T^{M_i}z_{k_i},z_{k_i}\right)\leq 2\varepsilon+\varepsilon_{k_{N+1}}\leq 3\varepsilon$(不妨设$\varepsilon_{k_{N+1}}<\varepsilon)$.\par
\end{proof}

\begin{definition}{齐性集}
    设$(X,d)$是一个紧致度量空间，$T:X\to X$ 同胚，$A\subseteq X$ 非空闭集，如果存在$S:X\to X$ 同胚，使得：
    \begin{enumerate}[(1)]
        \item $S\circ T=T\circ S$
        \item $S(A)=A$
        \item $S:A\to A$是极小系统
    \end{enumerate}
    则称$A$为\textbf{齐性集}.
\end{definition}

下面的引理说明如果一个齐性集对任意的“$\varepsilon$ 精度”都存在可以回到该精度的点，那么这个齐性集实际上是一个回复集.\par
\begin{lemma}{}{2.9}
    设$T:X\to X$ 是紧致度量空间$X$ 上的同胚，$A$ 为齐性集. 如果$\forall\, \varepsilon>0,\exists\, x,y\in A,n\in\mb{N}_{>0}$ ，使得$d \left(T^ny,x\right)<\varepsilon$，则$A$是回复集.
\end{lemma}
\begin{proof}
    由于$A$ 是齐性集，存在同胚映射$S:X\to X$满足： \par
    \begin{enumerate}[(1)]
        \item $S\circ T=T\circ S$
        \item $S(A)=A$
        \item $S:A\to A$是极小系统
    \end{enumerate}\par
    齐性集$A$是紧致度量空间中的闭集，因此$A$也是紧集，那么$\forall\,\varepsilon>0$ ，存在$A$ 的一个有限开覆盖$\{V_1,V_2,\cdots,V_m\}$，使得$\mathrm{diam}\, V_i\leq\varepsilon(1\leq i\leq m)$.\par
    由于$S:A\to A$是极小系统，则$\forall\, x\in A,\{\cdots,S^{-2}x,S^{-1}x,x,Sx,S^2x,\cdots\}$ 在$A$ 中稠密，因此$\{\cdots,S^{-2}x,S^{-1}x,x,Sx,S^2x,\cdots\}\cap V_i\neq\varnothing(\forall\, 1\leq i\leq m)$.\par
    $\Rightarrow \exists\, l\in\mb{Z},S^lx\in V_i\Rightarrow x\in S^lV_i$.那么$\left\{S^lV_i,l\in\mb{Z}\right\}$是$A$ 的一个开覆盖.\par
    由于$A$是紧集，存在有限子覆盖$S^{-l_i}V_i,S^{-(l_i-1)}V_i,\cdots,S^{l_i}V_i(\forall\, 1\leq i\leq m)$.\par
    令$L=\max_{1\leq i\leq m}{l_i}$，则$\left\{S^{-L}V_i,S^{-(L-1)}V_i,\cdots,V_i,\cdots,S^{L}V_i\right\}$是$A$ 的一个开覆盖.\par
    \begin{claim}{$\forall\, x,y\in A,\exists\, l\in\{-L,-(L-1),\cdots,L\}$ ，使得$d\left(S^ly,x\right)<\varepsilon$.}
        $\left\{V_i\right\}_{1\leq i\leq m}$ 是$A$ 的开覆盖$\Rightarrow \exists\, 1\leq i\leq m,\st x\in V_i$.\par
        又$\left\{S^{-L}V_i,S^{-(L-1)}V_i,\cdots,V_i,\cdots,S^{L}V_i\right\}$是$A$ 的开覆盖$\Rightarrow \exists\, k\in\{-L,\cdots,L\},\st y\in S^{k}V_i$.\par
        $\Rightarrow S^{-k}y,x\in V_i\Rightarrow d \left(S^{-k}y,x\right)<\varepsilon$，令$l=-k$即可.
    \end{claim}\par
    由引理条件得，$\forall\,\delta>0,\exists\, x_0,y_0\in A,n_0\in\mb{N}_{>0}$ ，使得$d \left(T^{n_0}y_0,x_0\right)<\delta$.\par
    由断言1，$\forall\, x\in A,\exists\, l\in\{-L,\cdots,L\}$ ，使得$d \left(S^lx_0,x\right)<\varepsilon$.\par
    $A$是紧集，因此同胚$S,T$ 一致连续，故$\forall\,\varepsilon>0,\exists\,\delta>0$ ，使得$d(x,y)<\delta\Rightarrow d \left(Tx,Ty\right)<\varepsilon,d \left(S^lx,S^ly\right)<\varepsilon(\forall\, -L\leq l\leq L)$.\par
    $\Rightarrow d \left(S^lT^{n_0}y_0,S^lx_0\right)\leq\varepsilon\xLongrightarrow{\text{可交换}} d \left(T^{n_0}\big(S^ly_0\big),S^lx_0\right)\leq\varepsilon\xLongrightarrow{\text{三角不等式}} d\left(T^{n_0}\big(S^ly_0\big),x\right)\leq2\varepsilon$.\par
    令$y=S^ly_0,n=n_0$即可.
\end{proof}

\begin{lemma}{}{2.10}
    设$(X,d)$是紧致度量空间，$F_i$是$X$ 中的闭集$(i\in\mb{N})$ ，如果$X=\bigcup_{i\in\mb{N}}{F_i}$ ，那么$\exists\, i_0\in\mb{N}$ ，使得$F_{i_0}$ 有内点.
\end{lemma}
\begin{proof}
    反证法.假设$\forall\, i\in\mb{N}$ ，$F_i$ 没有内点.\par
    对于$i=1$ ，存在$x_1\in X,\st x_1\notin F_1\Rightarrow \delta_1>0$ ，使得$\overline{B(x_1,\delta_1)}\cap F_1=\varnothing$.\par
    对于$i=2$ ，$B(x_1,\delta_1)\not\subset F_2$，否则$F_2$有内点，则$\exists\, x_2\in B(x_1,\delta_1),x_2\notin F_2$\par
    $\Rightarrow \exists\,\delta_2>0$，使得$B(x_2,\delta_2)\subset B(x_1,\delta_2)$且$\overline{B(x_2,\delta_2)}\cap F_2=\varnothing$.\par
    对于$\forall\, i\in\mb{N}$，存在$x_i\in B(x_{i-1},\delta_{i-1}),x_i\notin F_i\Rightarrow \exists\,\delta_i>0$，使得$B(x_i,\delta_i)\subset B(x_{i-1},\delta_{i-1})$，$\overline{B(x_i,\delta_i)}\cap F_i=\varnothing$ .\par
    设$B_0=\bigcap_{i\in\mb{N}}{\overline{B(x_i,\delta_i)}}$，那么$B_0\neq\varnothing$，设$x_0\in B_0$且$x_0\in F_j(j\in\mb{N})$.\par
    另一方面，$x_0\in\overline{B(x_j,\delta_j)}$，而$\overline{B(x_j,\delta_j)}\cap F_j=\varnothing$，矛盾.\par
    因此$\exists\, i_0\in\mb{N}$，$F_{i_0}$有内点.
\end{proof}
\begin{remark}
    上述引理的一个等价形式是Baire纲定理，即完备度量空间是第二纲集.
\end{remark}

\begin{lemma}
    设$(X,d)$是紧致度量空间，$T:X\to X$ 同胚，$A\subseteq X$ 闭集，如果$A$ 是齐性集，又是回复集，那么$A$一定有回复点，即$\exists\, x\in X,x\in\omega(x)$.
\end{lemma}
\begin{proof}
    定义$\setjot[-6]\begin{aligned}[t]
        F:A&\to \mb{R}_{>0} \\
        x&\mapsto\inf_{n\leq 1}{d \left(x,T^nx\right)}
    \end{aligned}$，则$A$中存在回复点$\iff\exists\, x\in A,F(x)=0$.\par
    反证法，假设$\forall\, x\in A,F(x)\neq 0$.\par
    令$S_n=\left\{x\in A\left|F(x)\geq\frac{1}{n}\right.\right\}(\forall\, n\in\mb{N})$，则$A=\bigcup_{n\in\mb{N}}{S_n}$.\par
    现证$S_n$是闭集.\par
    $S_n^c=\left\{x\in A\left|F(x)<\frac{1}{n}\right.\right\}$.则$\forall\, x\in S_n^c,\exists\, N\in\mb{N}$ ，使得$d \left(x,T^Nx\right)<\frac{1}{n}$.\par
    由$T$的连续性，故存在$x$的邻域$U_x$，使得$\forall\, y\in U_x,d \left(x,T^Nx\right)<\frac{1}{n}$.\par
    $\Rightarrow U_x\subseteq S_n^c\Rightarrow S_n^c$为开集，$S_n$为闭集.\par
    由引理\ref{lem:2.10}，$\exists\, n_0\in\mb{N}$，$S_{n_0}$包含内点.\par
    不妨设$B(x_0,r_0)\subset S_{n_0}$.\par
    由于$A$是齐性集，存在同胚映射$S:X\to X$满足： \par
    \begin{enumerate}[leftmargin=4em,label=(\arabic*)]
        \item $S\circ T=T\circ S$
        \item $S(A)=A$
        \item $S:A\to A$是极小系统
    \end{enumerate}\par
    $B(x_0,r_0)\subseteq S_{n_0}\subseteq A$，$S:A\to A$是极小系统$\Rightarrow \left\{S^l \left(B(x_0,r_0)\right)\Big| l\in\mb{Z}\right\}$是$A$的一个开覆盖.\par
    由$A$是闭集可知$A$是紧集，因此$\exists\, L\in\mb{N}$，使得$S^{-L}\left(B(x_0,r_0)\right)\cup S^{-(L-1)}\left(B(x_0,r_0)\right)\cup\cdots\cup S^{L}\left(B(x_0,r_0)\right)=A$.\par
    由于$S^{-L},\cdots,S^L$是紧集$X$上的连续函数，故一致连续$\Rightarrow\forall\,\varepsilon>0,\exists\,\delta>0$，使得$d(x,y)<\delta$，有$d \left(S^ix,S^iy\right)<\varepsilon(\forall\, i\in\{-L,\cdots,L\})$.\par
    令$\varepsilon=\frac{1}{n_0}$，$\exists\,\delta_0>0$，使得$d(x,y)<\delta_0\Rightarrow d \left(S^ix,S^iy\right)<\varepsilon=\frac{1}{n_0}\;(\forall\, i\in\{-L,\cdots,L\})$.\par
    \begin{claim}{$\forall\,y\in A,F(y)\geq \delta_0$}
        反证法，假设$\exists\, y_0\in A$，使得$F(y_0)<\delta_0\Rightarrow \exists\, N\in\mb{N},d \left(y_0,T^Ny_0\right)<\delta_0$.\par
        由于$y_0\in A=S^{-L}\left(B(x_0,r_0)\right)\cup S^{-(L-1)}\left(B(x_0,r_0)\right)\cup\cdots\cup S^{L}\left(B(x_0,r_0)\right)$，$\exists\, j_0\in\{-L,\cdots,L\}$，使得$y_0\in S^{j_0}\left(B(x_0,r_0)\right)$.\par
        $\Rightarrow S^{-j_0}(y_0)\in B(x_0,r_0)$，其中$-j_0\in \{-L,\cdots,L\}$.\par
        由于$d \left(y_0,T^Ny_0\right)<\delta_0\Rightarrow d \left(S^{-j_0}(y_0),S^{-j_0}\left(T^Ny_0\right)\right)<\varepsilon=\frac{1}{n_0}$.\par
        $\Rightarrow F \left(S^{-j_0}(y_0)\right)<\frac{1}{n_0}$，与$S^{-j_0}(y_0)\in B(x_0,r_0)\subseteq S_{n_0}$矛盾.
    \end{claim}\par
    因为$A$是回复集，由引理\ref{lem:2.8}知$\forall\, \varepsilon>0,\exists\, z\in A,n\in\mb{N}$ ，使得$d \left(z,T^nz\right)<\varepsilon$.\par
    取$\varepsilon=\frac{\delta_0}{2}\Rightarrow \exists\, z\in A,\st F(z)<\frac{\delta_0}{2}$，与断言1矛盾.
\end{proof}

最后，我们可以证明多重Birkhoff回复定理.
\begin{theorem*}{多重Birkhoff回复定理}
    设$(X,d)$ 是紧致度量空间，$T:X\to X$ 同胚，$l\in\mb{N}$，那么$\exists\, x\in X$，以及自然数子列$n_1<n_2<\cdots<n_k<\cdots$，使得$T^{n_i}x\to x,T^{2n_i}x\to x,\cdots,T^{ln_i}x\to x$.
\end{theorem*}
\begin{proof}
    对$l$应用数学归纳法.\par
    当$l=1$时，由Birkhoff回复定理知成立.\par
    考虑$\forall\, l\in\mb{N}$，假设$T:X\to X$是极小系统.\par
    则$Y=\underbrace{X\times\cdots\times X}_{l\,\text{个}\,X}$为紧致度量空间，$A=\left\{(x,x,\cdots,x)|x\in X\right\}$是闭集.\par
    定义映射$\setjot[-4]\begin{aligned}[t]
        \Phi:Y &\to Y \\
        (x_1,x_2,\cdots,x_l) &\mapsto \left(Tx_1,T^2x_2,\cdots,T^lx_l\right)
    \end{aligned}$，$\setjot[-4]\begin{aligned}[t]
        \Psi:Y &\to Y \\
        (x_1,x_2,\cdots,x_l) &\mapsto \left(Tx_1,Tx_3,\cdots,Tx_l\right)
    \end{aligned}$\par
    则$\Phi,\Psi$都是$Y$上的同胚，且$\Phi\circ\Psi=\Psi\circ\Phi,\Psi(A)=A$，由$T:X\to X$是极小系统可知$\Psi:A\to A$是极小系统.\par
    因此$A$是$X$的齐性集.\par
    由归纳假设，$\exists\, y\in Y$，以及一个自然数序列$n_1<n_2<\cdots<n_k<\cdots$，使得$T^{n_i}y\to y,\cdots,T^{(l-1)n_i}y\to y\;(i\to\infty)$.\par
    $\setjot[-2]\begin{aligned}[t]
        \Rightarrow &\left(T^{n_i}\left(T^{-n_i}y\right),T^{2n_i}\left(T^{-n_i}y\right),\cdots,T^{ln_i}\left(T^{-n_i}y\right)\right)\\
        &=\left(y,T^{n_i}y,\cdots,T^{(l-1)n_i}y\right)\to(y,y,\cdots,y)\in A
    \end{aligned}$\par
    则$\forall\,\varepsilon>0,\exists\,(y,y,\cdots,y)\in A,\left(T^{-n_i}y,\cdots,T^{-n_i}y\right)\in A$，以及$n_i\in\mb{N}$，使得$$d \left((y,\cdots,y),\Phi^{n_i}\left(T^{-n_i}y,\cdots,T^{-n_i}y\right)\right)<\varepsilon$$ \par
    由引理\ref{lem:2.9}，$A$是回复集.\par
    又$A$是齐性集，因此$A$中存在回复点$x\in \omega(x)$.\par
    即存在一个自然数序列$m_1<m_2<\cdots<m_k<\cdots$，使得$\Phi^{m_i}(\underbrace{(x,x,\cdots,x)}_{l\,\text{个}})\to\underbrace{(x,x,\cdots,x)}_{l\,\text{个}}$ \par
    $\Leftrightarrow T^{m_i}x\to x,T^{2m_i}x\to x,\cdots,T^{lm_i}\to x$
\end{proof}


\section{Van de Waerden定理}

下面使用Birkhoff多重回复定理证明数论中著名的Van de Waerden定理.\par
\begin{theorem}{Van de Waerden定理}
    设$\mb{N}=C_1\cup C_2\cup\cdots\cup C_l$是$\mb{N}$的一个有限分划$(l\in\mb{N})$，$\forall\, s\in\mb{N}$，存在$i\in\{1,\cdots,l\}$，使得$C_i$中有长度为$s$的等差数列.
\end{theorem}
\begin{proof}
    设$X=\{1,2,\cdots,l\}$，$Y=\prod_{i=-\infty}^\infty{X}$，$\setjot[-4]\begin{aligned}[t]
        T:Y&\to Y \\
        (\cdots,x_{-1},x_0,x_1,\cdots)&\mapsto (\cdots,x_0,x_1,x_2,\cdots)
    \end{aligned}$为左平移映射.\par
    $\forall\, x=(\cdots,x_{-1},x_0,x_1,\cdots),y=(\cdots,y_{-1},y_0,y_1,\cdots)\in Y$，定义$Y$上的度量为$d(x,y)=\sum_{i=-\infty}^\infty{\frac{|x_i-y_i|}{l^{|i|}}}$.\par
    \begin{claim}{$(Y,d)$是紧致度量空间，$T:Y\to Y$是连续映射.}
        只需证明$(Y,d)$是列紧的，即任何点列有收敛子列.\par
        $\forall\,\{y_n\}\subset Y,y_n=(\cdots,y_{-1}^n,y_0^n,y_1^n,\cdots)$.\par
        由于$X$有限，一定存在无穷子列$\{\prescript{0}{}y_n\}$，使得$\forall\, n,\prescript{0}{}y_0^n=y_0$.\par
        同理，存在无穷子列$\{\prescript{1}{}y_n\}\subseteq \{\prescript{0}{}y_n\}$，使得$\forall\, n,\prescript{1}{}y_{1}^n=y_{1}$.\par
        重复以上过程，依次可得子列$\{\prescript{1}{}y_n\}\supseteq\{\prescript{-1}{}y_n\}\supseteq\{\prescript{2}{}y_n\}\supseteq\{\prescript{-2}{}y_n\}\supseteq\cdots$，则$\exists\, y=(\cdots,y_{-1},y_0,y_1,\cdots)$，且可从上述点列中依次选取一个元素按序构成子列$\{\tilde{y}_n\}$，满足$\tilde{y}_i^n\to y_i$.\par
        现证$d(\tilde{y}_n,y)\to 0$.\par
        $d(\tilde{y}_n,y)=\sum_{i=-\infty}^{\infty}\frac{\left|\tilde{y}_i^n-y_i\right|}{l^{|l_i|}}=\sum_{i=-L}^L{\frac{\left|\tilde{y}_i^n-y_i\right|}{l^{|l_i|}}}+\sum_{|i|\geq L+1}{\frac{\left|\tilde{y}_i^n-y_i\right|}{l^{|l_i|}}}$.\par
        $\forall\,\varepsilon>0$,由于$\left|\tilde{y}_i^n-y_i\right|\leq l$，故$\exists\, L\in\mb{N}$，使得$\sum_{|i|\geq L+1}{\frac{\left|\tilde{y}_i^n-y_i\right|}{l^{|l_i|}}}<\varepsilon$.\par
        固定$L$，取$N=2L+1$，则当$n>N$时，$\forall\, -L\leq i\leq L,\tilde{y}_i^n=y_i\Rightarrow \sum_{i=-L}^L{\frac{\left|\tilde{y}_i^n-y_i\right|}{l^{|l_i|}}}=0$.\par
        故$\forall\, n>N$，有$d(\tilde{y}_n,y)<\varepsilon$，因此$d(\tilde{y}_n,y)\to 0$.
    \end{claim}\par
    \begin{claim}{$T:Y\to Y$是同胚}
        $\forall\, x=(\cdots,x_{-1},x_0,x_1,\cdots),y=(\cdots,y_{-1},y_0,y_1,\cdots),d \left(Tx,Ty\right)=\sum_{i=-\infty}^\infty{\frac{|x_{i+1}-y_{i+1}|}{l^{|i|}}}\leq \left(\sum_{i=-\infty}^\infty{\frac{|x_{i}-y_{i}|}{l^{|i|}}}\right)l=ld(x,y)$.\par
        类似地，$d \left(T^{-1}x,T^{-1}y\right)\leq ld(x,y)$.\par
        因此$T$是同胚.
    \end{claim}\par
    定义$p=(\cdots,p_{-1},p_0,p_1,\cdots)$.$\forall\,i\in\mb{Z}$，如果$|i|\in C_j(j\in\{1,2,\cdots,l\})$，那么$p_i=j$.\par
    有$p\in Y$，令$X_0=\overline{\left\{T^kp:k\in\mb{Z}\right\}}$，则$T:X_0\to X_0$是同胚，$(X_0,d)$是紧致度量空间.\par
    由Birkhoff多重回复定理，$\forall\, s\in\mb{N},\exists\,y\in X_0$以及一个自然数序列$n_1<n_2<\cdots<n_k<\cdots$，使得$T^{n_k}y\to y,T^{2n_k}y\to y,\cdots,T^{sn_k}y\to y(k\to\infty)$.\par
    由度量$d$的定义，如果$d(a,b)<1$，那么$a_0=b_0$，否则$d(a,b)\geq\frac{|a_0-b_0|}{l^0}\geq 1$.\par
    所以存在整数序列$\{m_k\}$，使得$T^{m_k}p\to y(k\to\infty)$.\par
    故$\exists\,m\in\mb{Z}$，使得$d \left(T^nT^mp,T^mp\right)<1,d(T^{2n}(T^mp),T^mp)\leq 1\cdots,T^{T^{sn}\left(T^mp\right),T^mp}\leq 1$.\par
    $\Rightarrow T^{m+n}p,T^{m+2n}p,\cdots,T^{m+sn}p$与$T^mp$在第0个位置相同.\par
    $T^kp=\left(\cdots,p_{k-1},p_k,p_{k+1},\cdots\right)$\par
    $\Rightarrow p_{m+n}=p_{m+2n}=\cdots=p_{m+sn}=p_m\Rightarrow \exists\,j\in\{1,2,\cdots,l\},|m+in|\in C_j\,(0\leq i\leq s)$.\par
    讨论$m+in$的符号，得$C_j$中有长度至少$\geq\frac{s}{2}$的等差数列
\end{proof}






